\section{Introduction}
\label{sec:introduction}

Automated theorem proving has made substantial progress through systems such
as AlphaProof~\citep{alphageometry2024} and
LeanDojo~\citep{yang2023leandojo}, yet a complementary capability has
received less attention: \emph{automated conjecture refutation}.
A system that can quickly disprove false conjectures serves two roles in a
research pipeline---it filters low-quality outputs from conjecture generators
and provides negative evidence that guides mathematical intuition.

The dominant paradigm in AI-for-mathematics treats the problem as proof
generation: given a statement believed to be true, find a proof.
Counterexample search occupies a narrower niche despite its practical
importance.
QuickCheck-style random testing~\citep{claessen2000quickcheck} works well for
programmable properties but scales poorly to number-theoretic or algebraic
statements requiring structured sampling.
Large language models (LLMs) can propose candidate counterexamples, but
prior work stops short of a principled multi-strategy framework with adaptive
resource allocation and formal symbolic verification.

We make the following contributions:

\begin{enumerate}[noitemsep]
  \item \textbf{ConjLean pipeline} (\S\ref{sec:system}): an end-to-end system
    that generates conjectures via a temperature-0.8 LLM, applies a SymPy
    symbolic pre-filter, formalizes survivors in Lean~4, and attempts proof
    via a four-layer escalating tactic search.

  \item \textbf{REFUTE pipeline} (\S\ref{sec:system}): a three-agent
    counterexample-search framework combining deterministic boundary sweep,
    structured random sampling, LLM-guided symbolic perturbation, and
    analogical transfer, coordinated by a UCB1 strategy-selection bandit.

  \item \textbf{A tiered benchmark} (\S\ref{sec:benchmark}): 92 conjectures
    across three difficulty tiers---synthetic-false with verified
    counterexamples (Tier~1), historical conjectures with known outcomes
    (Tier~2), and subtle or open cases (Tier~3)---spanning number theory,
    inequality, and combinatorics.

  \item \textbf{Empirical evaluation} (\S\ref{sec:results}): three runs
    comparing 7B, 32B, and a patched-verifier condition; scaling does not
    improve recall, a parser fix recovers 7 refutations, and LLM-generated
    candidate quality is identified as the binding constraint.
\end{enumerate}
